\subsection{Irreducible elements}

DEFINITION 6. --- An element x of an ordered group G is said to be irreducible if it is a minimal element in the set of strictly positive elements of G.

Let x be an irreducible element of the ordered group G; if y is a positive element of G then the element \(\inf(x, y)\) , if it exists, can only be equal to x or to 0. Thus in a lattice ordered group G, every positive y is either greater than or coprime to the irreducible element x; in particular, two distinct irreducible elements are coprime.

\begin{enumerate}
\def\labelenumi{(\Roman{enumi})}
\setcounter{enumi}{503}
\tightlist
\item
  An element \(p\) of \(\mathbf{A}\) is called irreducible if the ideal \((p)\) is an irreducible element of the ordered group \(\mathcal{P}^*\) ; this says that \(p\) is neither zero nor invertible, and that any element of \(\mathbf{A}\) which divides \(p\) is associate either to \(p\) or to 1. If \(\mathcal{P}^*\) is lattice ordered, then every \(a \in \mathbf{A}\) is either coprime to \(p\) or a multiple of \(p\) .
\end{enumerate}

Examples (DIV). --- 1) An integer \(p > 0\) is irreducible in \(\mathbf{Z}\) if and only if it is prime (I, p.~50).

\begin{enumerate}
\def\labelenumi{\arabic{enumi})}
\setcounter{enumi}{1}
\tightlist
\item
  A polynomial in one indeterminate over a field K is irreducible if and only if it is irreducible in the usual sense (IV, p.~13).
\end{enumerate}

PROPOSITION 14. --- For an element \(x > 0\) of an ordered group \(G\) to be irreducible it is sufficient that it satisfy the following property:

\begin{enumerate}
\def\labelenumi{(\Alph{enumi})}
\setcounter{enumi}{15}
\tightlist
\item
  The relations \(x \leqslant y + z\) , \(y \geqslant 0\) , \(z \geqslant 0\) imply \(x \leqslant y\) or \(x \leqslant z\) .
\end{enumerate}

This condition is necessary when G is lattice ordered.

If G is lattice ordered and x is irreducible, we have just seen that y is either greater than x or coprime to x; in the latter case Cor. 2 of VI, p.15 shows that z is greater than x. Conversely, suppose the condition is satisfied: if \(0 \leqslant y \leqslant x\) then it follows, by putting \(x = y + z\) ( \(z \geqslant 0\) ), that either \(x \leqslant y\) or \(x \leqslant z\) ; in the first case we have x = y; in the second we have \(x \leqslant x - y\) , so \(y \leqslant 0\) and so y = 0; this shows that x is indeed irreducible.

PROPOSITION 14 (DIV). --- For a nonzero element p of A to be irreducible it is sufficient that it not be a unit, and that it cannot divide a product of two elements of A without dividing one or other of them. This condition is necessary if \(P^{*}\) is lattice ordered.

Remark. --- Proposition 14 (DIV) can also be expressed as follows: if p is a non zero element of A such that the ideal \((p)\) is prime (I, p.~117, Def. 3) then p is irreducible; conversely, if \(P^{*}\) is lattice ordered and p is irreducible then the ideal \((p)\) is prime.

PROPOSITION 15. --- In an ordered group G, let \((p_{\iota})_{\iota \in I}\) be a family of pairwise distinct positive elements of G satisfying condition \((P)\) (and hence irreducible). Then the map

\[
(n _ {i}) _ {i \in I} \mapsto \sum_ {i \in I} n _ {i} p _ {i}
\]

is an isomorphism of the ordered group \(\mathbf{Z}^{(I)}\) , the direct sum of the ordered groups Z (VI, p.~7) onto the ordered subgroup of G generated by the \(p_{i}\) .

It is enough to show that the relation \(\sum_{\iota \in I} n_{\iota} p_{\iota} \geqslant 0\) is equivalent to \(n_{\iota} \geqslant 0\) for all \(\iota\) ,

for in particular the relation \(\sum_{\iota \in I} n_{\iota} p_{\iota} = 0\) will imply \(n_{\iota} = 0\) for all \(\iota\) , hence this will

show that the family \((p_{\iota})\) is linearly independent. Now let \(I'\) (resp. \(I''\) ) be the finite subset of \(I\) consisting of those \(\iota\) such that \(n_{\iota} > 0\) (resp. \(n_{\iota} < 0\) ); we have

\[
\sum_ {i \in I ^ {\prime}} n _ {i} p _ {i} \geqslant \sum_ {i \in I ^ {\prime \prime}} (- n _ {i}) p _ {i}.
\]

In particular, for \(\lambda \in I''\) , this implies that \(p_{\lambda} \leqslant \sum_{i \in I'} n_i p_i\) , and it follows by induction

from property (P) that we must have \(p_{\lambda} \leqslant p_{\iota}\) for some \(\iota \in I'\) ; since \(p_{\iota}\) is irreducible, this would imply \(p_{\lambda} = p_{\iota}\) , which is absurd. Hence \(I''\) is empty, which proves the proposition.

THEOREM 2. --- Let G be a filtered group. Then the following properties are equivalent :

\begin{enumerate}
\def\labelenumi{\alph{enumi})}
\item
  G is isomorphic to an ordered group of the form \(\mathbf{Z}^{(I)}\) .
\item
  G is lattice ordered and satisfies the following condition:
\end{enumerate}

(MIN) Every nonempty set of positive elements of \(\mathbf{G}\) has a minimal element.

\begin{enumerate}
\def\labelenumi{\alph{enumi})}
\setcounter{enumi}{2}
\item
  G satisfies condition (MIN) and every irreducible element of G has property (P).
\item
  G is generated by its irreducible elements, and every irreducible element of G has property (P).
\end{enumerate}

Let us show first that \(a\) implies \(b\) . The group \(\mathbf{Z}^{(I)}\) is lattice ordered, as the direct sum of totally ordered groups. On the other hand let E be a nonempty set of positive elements of \(\mathbf{Z}^{(I)}\) and let \(x = \sum_{i} n_i e_i\) be an element of E (where

\((e_{i})\) denotes the natural basis of \(\mathbf{Z}^{(I)}\) ; there are a finite number \(\prod_{i}(n_{i} + 1)\) of

elements \(y\) of \(\mathbf{Z}^{(I)}\) such that \(0 \leqslant y \leqslant x\) , so the set \(F\) of elements of \(E\) which are less than or equal to \(x\) is a fortiori finite; since it is nonempty, it contains a minimal element (Set Theory, III, p.~170, Cor. 2), which is clearly a minimal element of \(E\) .

It is clear that \(b)\) implies \(c)\) , by Prop. 14. Let us show that \(c)\) implies \(d)\) . Since \(G\) is filtered, it is enough (VI, p.~4, Prop. 4) to check that the set \(F\) of positive elements of \(G\) which are sums of irreducible elements is equal to \(G_{+} - \{0\}\) . If this were not true, it would follow from (MIN) that the complement of \(F\) in \(G_{+} - \{0\}\) would have a minimal element \(a\) ; by definition \(a\) is not irreducible, so is the sum of two strictly positive elements \(x\) and \(y\) ; since \(x < a\) and \(y < a\) , these elements belong to \(F\) , and so are sums of irreducible elements, and it follows that so is \(a\) , which is a contradiction. Finally, \(d)\) implies \(a)\) by Prop. 15.

We will apply Th. 2 to the theory of divisibility in principal ideal domains (VII, p.~4) and in unique factorisation domains (AC, VII, § 3), as well as to the study of ideals in a Dedekind ring (AC, VII, § 2).

\section{ORDERED FIELDS}

\subsection{Ordered rings}

DEFINITION 1. --- Given a commutative ring A, we say that an ordering on A is compatible with the ring structure of A if it is compatible with the additive group structure of A, and if it satisfies the following axiom:

(OR) The relations \(x \geqslant 0\) and \(y \geqslant 0\) imply \(xy \geqslant 0\) .

The ring A, together with such an ordering, is called an ordered ring.

Examples. --- 1) The rings Q and Z, with the usual orderings, are ordered rings. 2) A product of ordered rings, equipped with the product ordering, is an ordered ring. In particular the ring \(A^{E}\) of mappings from a set E to an ordered ring A is an ordered ring.

\begin{enumerate}
\def\labelenumi{\arabic{enumi})}
\setcounter{enumi}{2}
\tightlist
\item
  A subring of an ordered ring, with the induced ordering, is an ordered ring.
\end{enumerate}

In an ordered ring, the relations \(x \geqslant y\) and \(z \geqslant 0\) imply \(xz \geqslant yz\) . Indeed these inequalities are equivalent to \(x - y \geqslant 0\) , \(z \geqslant 0\) and \((x - y)z \geqslant 0\) respectively.

Analogously we can show that the relations \(x \leqslant 0\) and \(y \geqslant 0\) (resp. \(y \leqslant 0\) ) imply \(xy \leqslant 0\) (resp. \(xy \geqslant 0\) ). These results are often invoked under the name of sign rules (two elements are said to have the same sign if they are both \(\geqslant 0\) or both \(\leqslant 0\) ). They imply that, if A is a totally ordered ring, then every square is positive, and in particular that every idempotent (for example the unit element) is positive.

Example. --- There is only one totally ordered ring structure on Z: indeed 1 \textgreater{} 0, whence n \textgreater{} 0 for every natural number \(n \neq 0\) , by induction. In contrast there exist ordered ring structures on Z which are not totally ordered (see below).

Let P be the set of positive elements of an ordered ring A. It is known (VI, p.~3, Prop. 3) that P determines the ordering on A. To say that A is an ordered ring is equivalent to saying that P satisfies the following properties:

\[
\begin{array}{l l} \left(\mathrm{AP} _ {\mathrm{I}}\right) & \mathrm{P} + \mathrm{P} \subset \mathrm{P} \\ \left(\mathrm{AP} _ {\mathrm{II}}\right) & \mathrm{PP} \subset \mathrm{P} \\ \left(\mathrm{AP} _ {\mathrm{III}}\right) & \mathrm{P} \cap (- \mathrm{P}) = \{0 \}. \end{array}
\]

Indeed \((\mathrm{AP}_{\mathrm{I}})\) and \((\mathrm{AP}_{\mathrm{III}})\) state that the additive group of A is an ordered group (VI, p.~3, Prop. 3), while \((\mathrm{AP}_{\mathrm{II}})\) is a translation of (OR).

Recall that the following condition is necessary and sufficient for the order relation on A to be total :

\[
\left(\mathrm{AP} _ {\mathrm{IV}}\right) \quad \mathrm{P} \cup (- \mathrm{P}) = \mathrm{A}.
\]

Example. --- In Z, if we take P to be the set of positive (in the usual sense) even integers, we get a ring which is not totally ordered.

Recall also that, in a totally ordered abelian group, the relation \(n \cdot x = 0\) (for a natural number \(n \neq 0\) ) implies \(x = 0\) (VI, p.~4); this gives us the following result.

PROPOSITION 1. --- A totally ordered ring is torsion free as a Z-module (II, p.~313).

\subsection{Ordered fields}

DEFINITION 2. --- A commutative field, equipped with a total ordering, is called an ordered field if its ordering and its ring structure are compatible.

We restrict ourselves to total order relations on fields because the others are very « pathological » (VI, p.~38, Ex. 6).

Examples. --- 1) The field \(\mathbf{Q}\) of rational numbers is an ordered field.

\begin{enumerate}
\def\labelenumi{\arabic{enumi})}
\setcounter{enumi}{1}
\item
  A subfield of an ordered field, with the induced ordering, is an ordered field.
\item
  * The field of real numbers is an ordered field. *
\end{enumerate}

Let \(\mathbf{K}\) be an ordered field. For all \(x\in \mathbf{K}\) we put

\[
\begin{array}{l l} \operatorname{sgn} (x) = 1 & \text {if} \quad x > 0, \\ \operatorname{sgn} (x) = - 1 & \text {if} \quad x <   0, \\ \operatorname{sgn} (x) = 0 & \text {if} \quad x = 0. \end{array}
\]

Then we have \(\operatorname{sgn}(xy) = \operatorname{sgn}(x)\operatorname{sgn}(y)\) ; we call \(\operatorname{sgn}(x)\) the sign of \(x\) . The map \(x \mapsto \operatorname{sgn}(x)\) from \(\mathbf{K}^*\) to the multiplicative group \(\{-1, +1\}\) is a surjective homomorphism whose kernel, the set of strictly positive elements of \(\mathbf{K}\) , is a subgroup of \(\mathbf{K}^*\) of index 2.

Conversely, if K is a commutative field and \(s: K^{*} \to \{-1, +1\}\) is a surjective homomorphism whose kernel is closed under addition, then s is the sign map for a unique ordered field structure, where the set of strictly positive elements is the kernel of s.

For all \(x\) and \(y\) in \(\mathbf{K}\) we have \(x = \operatorname{sgn}(x)|x|\) and \(|xy| = |x||y|\) .

On the other hand every ordered field is of characteristic zero (Prop. 1).

PROPOSITION 2. --- Let A be a totally ordered integral domain, and let K be its field of fractions. Then there exists one and only one ordering on K which restricts to the given ordering on A and makes K an ordered field.

Every \(x \in K\) can be expressed in the form \(x = ab^{-1}\) , with a and b in A and \(b \neq 0\) . If x is positive, then a and b have the same sign, and conversely. Thus we see that, if there exists an ordering on K satisfying the prescribed conditions, then it is unique, and the set P of positive elements is identical to the set of \(ab^{-1}\) , where a and b are elements of A of the same sign, and \(b \neq 0\) . It remains to show that P satisfies conditions \((\mathrm{AP}_{\mathrm{I}}), (\mathrm{AP}_{\mathrm{II}}), (\mathrm{AP}_{\mathrm{III}})\) and \((\mathrm{AP}_{\mathrm{IV}})\) . This is obvious for \((\mathrm{AP}_{\mathrm{II}})\) and \((\mathrm{AP}_{\mathrm{IV}})\) . For \((\mathrm{AP}_{\mathrm{I}})\) , consider \(ab^{-1} + cd^{-1}\) , where we may assume that \(a, b, c\) and \(d\) are positive; this sum is \((ad + bc)(bd)^{-1}\) , and \(ad + bc\) and \(bd\) are positive.

To show \((\mathrm{AP}_{\mathrm{III}})\) , consider an identity of the form \(ab^{-1} = -cd^{-1}\) , so that \(ad + bc = 0\) . If we assume that a and b have the same sign and that c and d have the same sign, then the sign rules show that ad and bc have the same sign; whence ad = bc = 0, so a = c = 0; hence P does indeed satisfy \((\mathrm{AP}_{\mathrm{III}})\) .

Example. --- Since Z admits only one totally ordered ring structure (VI, p.~19, example), the field Q admits only one ordering which makes it an ordered field: this is the usual ordering.

\subsection{Extensions of ordered fields}

DEFINITION 3. --- Let K be an ordered field. An ordered extension of K is a pair (E, u), where E is an ordered field and u is an increasing homomorphism from K to E.

Let K be a field, let E be an ordered field and let \(u: K \rightarrow E\) be a homomorphism. The relation

\[
x \leqslant y \quad \text { if } \quad u (x) \leqslant u (y)
\]

is a total order relation on K which gives it an ordered field structure, said to be induced by that of E. If K and E are ordered fields, then a homomorphism \(u: K \rightarrow E\) is increasing if and only if the ordered field structure of K is induced by that of E. We will usually identify K with its image in E under u.

Examples. --- 1) Every ordered field K is an ordered extension of Q. Indeed K is an extension of Q, since it is of characteristic zero, and on the other hand Q can only be ordered in one way, as we have just seen.

\begin{enumerate}
\def\labelenumi{\arabic{enumi})}
\setcounter{enumi}{1}
\tightlist
\item
  Let K be an ordered field, and let K(X) be the field of rational functions in one indeterminate over K. Let us define an ordering on the polynomial ring K{[}X{]} by taking the positive elements to be 0 and those polynomials whose leading coefficient is positive. In this way we obtain a totally ordered ring whose ordering extends that of K. By applying Prop. 2 we give K(X) the structure of an ordered extension of K. * For K = R it can be shown that the order relation defined on K(X) in this way is that of growth near +∞ (cf.~VI, p.~24, Prop. 4). *
\end{enumerate}

THEOREM 1. --- For an extension E of K to admit the structure of an ordered extension of K, the following condition is necessary and sufficient :

(OE) The relation \(p_{1}x_{1}^{2} + \cdots + p_{n}x_{n}^{2} = 0\) implies

\[
p _ {1} x _ {1} = \dots = p _ {n} x _ {n} = 0
\]

for any finite sequence \((x_{i},p_{i})\) of pairs of elements \(x_{i}\) of E and positive elements \(p_i\) of K.

Condition (OE) is clearly equivalent to :

(OE') The element \(-1\) is not a sum of elements of the form \(px^2\) ( \(x \in \mathbf{E}\) , \(p \in \mathbf{K}\) , \(p \geqslant 0\) ).

Condition (OE) is necessary: if E is an ordered extension of K then the elements \(p_{i}x_{i}^{2}\) are positive in E, so zero if their sum is zero. On the other hand \(p_{i}x_{i}^{2}=0\) is equivalent to \(p_{i}x_{i}=0\) .

Conversely, suppose condition (OE) is satisfied, then we will define an ordering on E by constructing a subset P of E which satisfies conditions \((\mathrm{AP}_{\mathrm{I}})\) , \((\mathrm{AP}_{\mathrm{II}})\) , \((\mathrm{AP}_{\mathrm{III}})\) and \((\mathrm{AP}_{\mathrm{IV}})\) , and which contains the set \(K_{+}\) of positive elements of K. Such a subset P will certainly make E an ordered extension of K, for we will have \(K \cap P = K_{+}\) ; indeed, if P were to contain an element -a \textless{} 0 of K, then a would belong to \(P \cap (-P)\) , contradicting \((\mathrm{AP}_{\mathrm{III}})\) .

To define P, let us consider the set M of subsets of E which satisfy \((\mathrm{AP}_{\mathrm{I}})\) , \((\mathrm{AP}_{\mathrm{II}})\) and \((\mathrm{AP}_{\mathrm{III}})\) , and which contain the union of \(K_{+}\) and the set C of squares of elements of E. This set M is nonempty, for it contains the set \(P_{0}\) of elements of the form \(\sum_{i} p_{i} x_{i}^{2}\) (that \(P_{0}\) satisfies \((\mathrm{AP}_{\mathrm{III}})\) follows immediately

from (OE)).

Moreover M is inductive (Set Theory, III, p.~154, Def. 3). Thus there exists, by Th. 2 of Set Theory, III, p.~154, a maximal element in M, which it remains for us to prove satisfies \((\mathrm{AP}_{\mathrm{IV}})\) ; now this follows from the following lemma:

Lemma. --- Let \(P \in M\) and \(x \notin P\) ; then there exists \(P' \in M\) such that \(P \subset P'\) and \(-x \in P'\) .

Take \(P' = P - xP\) , and check that \(P'\) has the required properties. Since \(0 \in C \subset P\) , we have \(P \subset P'\) . Whence \(C \subset P'\) and \(K_{+} \subset P'\) . Since \(1 \in C \subset P\) we have \(-x \in P'\) . We have

\[
\mathrm{P} ^ {\prime} + \mathrm{P} ^ {\prime} = \mathrm{P} - x \mathrm{P} + \mathrm{P} - x \mathrm{P} = \mathrm{P} + \mathrm{P} - x (\mathrm{P} + \mathrm{P}) \subset \mathrm{P} - x \mathrm{P} = \mathrm{P} ^ {\prime},
\]

whence \((\mathbf{AP}_1)\) . We have

\[
\begin{array}{r l} \mathrm {P^ {\prime} P^ {\prime}} = & (\mathrm{P} - x \mathrm{P}) (\mathrm{P} - x \mathrm{P}) \subset \\ & \subset \mathrm{PP} + x ^ {2} \mathrm{PP} - x (\mathrm{PP} + \mathrm{PP}) \subset \mathrm{P} + \mathrm{CP} - x \mathrm{P} \subset \mathrm{P} - x \mathrm{P} = \mathrm {P^ {\prime}}, \end{array}
\]

whence \((\mathrm{AP}_{\mathrm{II}})\) . Finally, let us check \((\mathrm{AP}_{\mathrm{III}})\) : suppose given an identity of the form \(p - xq = -(r - xs)\) where p, q, r, s belong to P; we deduce from this the relation \(x(s + q) = p + r\) ; if \(s + q \neq 0\) we have

\[
x = (s + q) ^ {- 2} (s + q) (p + r) \in \mathrm{CPP} \subset \mathrm{P},
\]

contrary to the hypothesis; hence \(s + q = 0\) , whence \(p + r = 0\) ; since \(\mathbf{P}\) satisfies \((\mathbf{AP}_{\mathrm{III}})\) we deduce that \(s = q = r = p = 0\) , which completes the proof.

COROLLARY 1 (« Artin-Schreier Theorem »). --- A necessary and sufficient condition for there to exist an ordering on a commutative field E which makes it an ordered field, is that the relation \(x_{1}^{2} + \cdots + x_{n}^{2} = 0\) imply \(x_{1} = \cdots = x_{n} = 0\) .

The necessity is obvious. Conversely, the stated condition implies that E is of characteristic zero, hence an extension of Q; then condition (OE) is satisfied, and Th. 1 shows that there exists on E the structure of an ordered extension of Q, that is an ordered field structure.

There does not exist any ordered field structure on a field E in which -1 is a square, in particular on an algebraically closed field.

COROLLARY 2. --- Let E be an extension of K admitting the structure of an ordered extension of K. For an element \(x \in E\) to be positive under every such structure on E, it is necessary and sufficient that x be of the form \(\sum_{i} p_{i} x_{i}^{2}\) , where \(x_{i} \in E\) and the

\(p_i\) are positive elements of \(\mathbf{K}\) .

The condition is obviously sufficient; it is also necessary, for (in the notation of the proof of Th. 1), if \(x \notin \mathbf{P}_0\) there exists a maximal element \(\mathbf{P}\) of \(\mathfrak{M}\) such that \(x \notin \mathbf{P}\) ; then \(-x \in \mathbf{P}\) by the Lemma, and \(x\) is not positive under the ordering defined by \(\mathbf{P}\) , since \(x \neq 0\) .

\subsection{Algebraic extensions of ordered fields}

Let K be an ordered field, and f a polynomial in K{[}X{]}. We will say that f changes sign in K if there exist two elements a and b in K such that \(f(a)\) \(f(b) < 0\) ; then we say that f changes sign between a and b.

PROPOSITION 3. --- Let K be an ordered field and f an irreducible polynomial over K which changes sign between a and b in K. Then the extension \(E = K[X]/(f)\) of K admits the structure of an ordered extension.

We will argue by induction on the degree n of f.~For n = 1 the proof is trivial. Suppose the result holds for degrees \(\leqslant n - 1\) , and let us prove it for degree n by contradiction; by Th. 1 we are thus assuming a relation of the form

\(1 + \sum_{i}p_{i}f_{i}^{2}(\mathbf{X})\equiv 0 (\mathrm{mod} f(\mathbf{X}))\) , where \(f_{i}\in \mathbf{K}[\mathbf{X}], p_{i}\in \mathbf{K}\) and \(p_i\geqslant 0\)

Without loss of generality we may suppose that the \(f_{i}\) have degrees \(\leqslant n - 1\) (IV, p.~11, Cor). Then

\[
1 + \sum_ {i} p _ {i} f _ {i} ^ {2} (\mathbf {X}) = h (\mathbf {X}) f (\mathbf {X})
\]

where \(h \neq 0\) has degree at most n - 2. Replacing X by a and b in the above inequality, we see that \(h(a) f(a) > 0\) and \(h(b) f(b) > 0\) . Since f changes sign between a and b by hypothesis, we conclude that \(h(a) h(b) < 0\) . Then we have a similar inequality for one of the irreducible factors \(g(\mathbf{X})\) of \(h(\mathbf{X})\) : that is \(g(a)g(b) < 0\) . But \(1 + \sum_{i}p_{i}f_{i}^{2}(\mathbf{X})\equiv 0 (\mathrm{mod} g(\mathbf{X}))\) , which shows that the field \(\mathrm{K}[\mathrm{X}] / (g)\) cannot be an ordered extension of \(\mathbf{K}\) (Th. 1), contrary to the induction hypothesis.

Remark. --- There exist irreducible polynomials f over an ordered field K which do not change sign in K, but such that \(\mathrm{K}[\mathrm{X}]/(f)\) admits the structure of an ordered extension of K (cf.~VI, p.~43, Ex. 26, c)).

In order to apply the previous proposition we will need the following result:

PROPOSITION 4. --- Let K be an ordered field and let \(f \in \mathbf{K}[\mathbf{X}]\) . There exists an interval in K, in the complement of which \(f\) takes the same sign as its highest degree term.

We can immediately reduce ourselves to the case of a monic polynomial; then one can write \(f(x) = x^{n}(1 + a_{1}x^{-1} + \cdots + a_{n}x^{-n})\) for \(x \neq 0\) . Let

\[
\mathbf {M} = \sup \left(1, | a _ {1} | + \dots + | a _ {n} |\right).
\]

For \(|x| > M\) we have \(1 + a_{1}x^{-1} + \cdots + a_{n}x^{-n} > 0\) , which completes the proof of the proposition.

COROLLARY 1. --- Every extension of an ordered field of odd finite degree admits the structure of an ordered extension.

Such an extension, being monogenous (V, p.~40, Th. 1), is isomorphic to K{[}X{]}/(f), where f is an irreducible polynomial of odd degree. Then it is enough to show that f changes sign in K (Prop. 3), which follows immediately from Prop. 4.

COROLLARY 2. --- If a is a positive element of an ordered field K, then every splitting field E of the polynomial \(X^{2}-a\) admits the structure of an ordered extension of K.

The result is trivial if a is a square in K. Otherwise the polynomial \(f(\mathbf{X}) = \mathbf{X}^{2} - a\) is irreducible and changes sign, since \(f(0) < 0\) and \(f(x)\) has the same sign as \(x^{2}\) , so positive, for x in the complement of some interval of K. We can now complete the proof by applying Prop. 3.

Remark. --- When the ordered field K contains the « square roots » of a positive element \(a\) of K (roots of the polynomial \(\mathbf{X}^2 - a\) ) then the notation \(\sqrt{a}\) is generally reserved for the positive square root. If K does not contain the square roots \(b\) and \(-b\) of \(a\) in the field E, then the latter can be made an ordered extension of K in two ways, each induced from the other via the K-automorphism which sends \(b\) to \(-b\) ; the choice of one of these orderings determines \(\sqrt{a}\) : it is whichever of the elements \(b\) and \(-b\) is positive.

If \(\pmb{a}\) and \(a'\) are two positive elements of \(\mathbf{K}\) , whose square roots are in \(\mathbf{K}\) , then \(\sqrt{aa'} = \sqrt{a}\sqrt{a'}\) , which follows from the definition of \(\sqrt{a}\) and the sign rule.

\subsection{Maximal ordered fields}

DEFINITION 4. --- An ordered field K is maximal if every ordered algebraic extension of K is trivial.

Example. --- * We will see later (Gen.~Top., VIII, p.~1) that the field R of real numbers is a maximal ordered field. *

The existence of maximal ordered fields is a consequence of the following theorem:

THEOREM 2. --- Every ordered field K admits an ordered algebraic extension which is a maximal ordered field.

One can show that this ordered extension is unique up to K-isomorphism (VI, p.~40, Ex. 15).

Let \(\Omega\) be an algebraic closure of \(K\) , and let \(\mathfrak{N}\) be the set of pairs \((A, \omega)\) , where \(A\) is a sub- \(K\) -extension of \(\Omega\) , and \(\omega\) is an ordering on \(A\) making \(A\) an ordered extension of \(K\) . Order \(\mathfrak{N}\) by the relation « \(L\) is an ordered extension of \(M\) » between \(M\) and \(L\) . Equipped with this ordering, \(\mathfrak{N}\) is an inductive ordered set: indeed if \((L_{t})\) is a totally ordered family of elements of \(\mathfrak{N}\) , then the field \(L = \cup_{t} L_{t}\) , ordered by taking \(L_{+} = \cup_{t} (L_{t})_{+}\) is an upper bound for the

\(L_{i}\) . Then N has a maximal element, by Set Theory, III, p.~154, Th. 2, which completes the proof.

PROPOSITION 5. --- Let K be a maximal ordered field, and let f be a polynomial in K{[}X{]} which changes sign between two elements a and b of K (with a \textless{} b). Then f has a root x in K such that a \textless{} x \textless{} b.

At least one of the irreducible factors of f, say h, changes sign between a and b. Then the field \(\mathrm{K}[\mathrm{X}]/(h)\) admits the structure of an ordered extension of K (VI, p.~23, Prop. 3), and h has degree 1 (Def. 4). Since \(h(a)h(b)<0\) , the unique root x of h is such that a\textless x\textless b, since a polynomial function of degree 1 is monotonic.

PROPOSITION 6. --- Every positive element of a maximal ordered field K has a square root in K. Every polynomial of odd degree in K{[}X{]} has at least one root in K.

This follows immediately from Cor. 2 and 1 to Prop. 4 of VI, p.~24.

COROLLARY. --- On a maximal ordered field K there exists only one ordering compatible with the field structure.

Indeed the positive elements of \(\mathbf{K}\) are determined by its algebraic structure: they are the squares.

\subsection{Characterisation of maximal ordered fields. Euler-Lagrange Theorem}

The property expressed by Prop. 6 of VI, p.~25 characterises maximal ordered fields. More precisely :

THEOREM 3 (Euler-Lagrange). --- Let K be an ordered field. Then the following three properties are equivalent :

\begin{enumerate}
\def\labelenumi{\alph{enumi})}
\item
  The field \(\mathbf{K}(i)\) is algebraically closed (where \(i\) denotes a square root of \(-1\) ).
\item
  The ordered field \(\mathbf{K}\) is maximal.
\item
  Every positive element of \(\mathbf{K}\) is a square, and every polynomial of odd degree in \(\mathbf{K}[\mathbf{X}]\) has a root in \(\mathbf{K}\) .
\end{enumerate}

It is clear that \(a)\) implies \(b)\) : indeed \(K\) has only two algebraic extensions up to isomorphism, the field \(K\) itself and \(K(i)\) , which cannot be ordered since \(-1\) is a square.

The fact that \(b)\) implies \(c)\) is nothing other than Prop. 6 of VI, p.~25.

It remains for us to prove that \(c)\) implies \(a)\) . That will follow from the next two propositions.

PROPOSITION 7. --- Let K be an ordered field in which every positive element is a square. Then every element of K(i) is a square, and every polynomial of degree 2 over K(i) has a root in K(i).

Let us show first that the second assertion reduces to the first. One can put the second degree polynomial \(aX^{2} + bX + c\) ( \(a \neq 0\) ) in the following form, often called the canonical trinomial form:

\[
a \left(\left(X + (b / 2 a)\right) ^ {2} - \left(b ^ {2} - 4 a c\right) / 4 a ^ {2}\right).
\]

If \(d\) is a square root of \((b^2 - 4ac) / 4a^2\) , then \(d - (b / 2a)\) is a root of the quadratic polynomial under consideration.

Now we show that every element \(a + bi\) ( \(a \in \mathbf{K}\) , \(b \in \mathbf{K}\) ) is a square; we are looking for an element \(x + yi\) such that

\[
(x + y i) ^ {2} = a + b i;
\]

this translates into \(x^{2} - y^{2} = a\) and \(2xy = b\) . From this we deduce that

\[
(x ^ {2} + y ^ {2}) ^ {2} = a ^ {2} + b ^ {2}.
\]

Let \(c\) denote the positive square root of \(a^2 + b^2\) ; then \(c \geqslant |a|\) , \(c \geqslant |b|\) and \(x^2 + y^2 = c\) . Whence \(x^2 = (c + a)/2\) and \(y^2 = (c - a)/2\) . Since \(c \geqslant |a|\) these equations are soluble in K, and if \(x_0\) and \(y_0\) are two solutions then \(x_0^2 - y_0^2 = a\) and \(2x_0y_0 = \pm b\) . We obtain the desired square root by taking \(x = x_0\) and \(y = b/2x_0\) .

PROPOSITION 8. --- Let K be a commutative field (of arbitrary characteristic) and let \(\mathbf{K}'\) be a splitting field for the polynomial \(\mathbf{X}^2 + 1 \in \mathbf{K}[\mathbf{X}]\) (V, p.~21). Suppose :

\begin{enumerate}
\def\labelenumi{\alph{enumi})}
\item
  every polynomial in K{[}X{]} of odd degree has a root in K';
\item
  every polynomial in \(\mathbf{K}'[\mathbf{X}]\) of degree 2 has a root in \(\mathbf{K}'\) . Then \(\mathbf{K}'\) is algebraically closed.
\end{enumerate}

Note first that it is enough to prove that every non-constant polynomial in K{[}X{]} has a root in K': this is indeed clear if \(K' = K\) ; if \(K' \neq K\) then \([K': K] = 2\) ; let \(a \mapsto \bar{a}\) denote the unique K-automorphism of \(K'\) distinct from the identity map; if \(f \in K'[X]\) and if \(\bar{f}\) denotes the polynomial obtained by applying \(a \mapsto \bar{a}\) to the coefficients of f, then \(f\bar{f} \in K[X]\) ; if \(a \in K'\) is a root of \(f\bar{f}\) then a is either a root of f or of \(\bar{f}\) ; thus either a or \(\bar{a}\) is a root of f.

Thus let f be a polynomial over K of degree \(2^{n}p\) , p odd. We will proceed by induction on n, the property being true for n = 0 by hypothesis a). Let E be an extension of K in which f splits into linear factors:

\[
f (\mathbf {X}) = \prod_ {i} \left(\mathbf {X} - a _ {i}\right).
\]

Let \(b \in K\) ; put \(y_{ij} = a_i + a_j + ba_i a_j \in E\) and

\[
h (\mathbf {X}) = \prod_ {i <   j} \left(\mathbf {X} - y _ {i j}\right) \in \mathrm{E} [ \mathbf {X} ].
\]

The coefficients of this polynomial are symmetric functions in the \(a_i\) , with coefficients in \(\mathbf{K}\) ; it therefore belongs to \(\mathbf{K}[\mathbf{X}]\) (IV, p.~62, Th. 1); since it has degree \(2^n p(2^n p - 1)/2 = 2^{n-1}p'\) ( \(p'\) odd), it has a root \(y_{ij}\) in \(\mathbf{K}'\) by inductive hypothesis. If we note that this holds for all \(b \in \mathbf{K}\) , and that \(\mathbf{K}\) is an infinite field (indeed a finite field, which has monogenous extensions of arbitrarily large odd degree (V, p.~94, Prop. 3), cannot satisfy \(a\) )), then we can deduce the existence of at least one pair \((i,j)\) such that

\[
a _ {i} + a _ {j} + b a _ {i} a _ {j} \in \mathbf {K} ^ {\prime} \quad \text { and } \quad a _ {i} + a _ {j} + b ^ {\prime} a _ {i} a _ {j} \in \mathbf {K} ^ {\prime},
\]

with \(b \neq b'\) . Then \(a_i + a_j\) and \(a_i a_j\) are elements of \(K'\) , hence so are \(a_i\) and \(a_j\) , since they are the roots of the quadratic equation

\[
x ^ {2} - \left(a _ {i} + a _ {j}\right) x + a _ {i} a _ {j} = 0.\tag{Q.E.D.}
\]

For a generalisation and an alternative proof of Prop. 8, based on Galois theory, see VI, p.~46, Ex. 33.

Let K be an ordered field and let \(\mathbf{K}^{\prime} = \mathbf{K}(i)\) ; for every element \(z = a + bi\) of \(K^{\prime}\) , the norm \(z\overline{z} = a^{2} + b^{2}\) of z relative to K (III, p.~544, example 1) is a positive element of K, which vanishes only for z = 0. If every positive element in K is a square (in particular if K is a maximal ordered field), then the positive square root of the norm \(z\overline{z}\) is called the absolute value of z, and is written \(|z|\) . Since \(|zz^{\prime}|^{2} = |z|^{2} |z^{\prime}|^{2}\) we have \(|zz^{\prime}| = |z| \cdot |z^{\prime}|\) .

Moreover, the triangle inequality

\[
\left| z + z ^ {\prime} \right| \leqslant \left| z \right| + \left| z ^ {\prime} \right|
\]

holds for every pair of elements \(z, z'\) of \(\mathbf{K}'\) . Indeed, if \(z = a + bi\) and \(z' = a' + b'i\) , then this inequality is equivalent to

\[
(a + a ^ {\prime}) ^ {2} + (b + b ^ {\prime}) ^ {2} \leqslant a ^ {2} + b ^ {2} + a ^ {\prime 2} + b ^ {\prime 2} + 2 \sqrt {(a ^ {2} + b ^ {2}) (a ^ {\prime 2} + b ^ {\prime 2})}
\]

and hence also to

\[
(a a ^ {\prime} + b b ^ {\prime}) ^ {2} \leqslant (a ^ {2} + b ^ {2}) (a ^ {\prime 2} + b ^ {\prime 2})
\]

which can be written \((ab' - ba')^2 \geqslant 0\) .

Th. 3 enables us to determine all the irreducible polynomials over a maximal ordered field :

PROPOSITION 9. --- If K is a maximal ordered field, then the only irreducible polynomials in K{[}X{]} are the first degree polynomials, and the second degree polynomials \(aX^{2} + bX + c\) such that \(b^{2} - 4ac < 0\) .

Since \(\mathbf{K}(i)\) is algebraically closed, every algebraic extension of K, and hence also every irreducible polynomial over K, has degree 1 or 2. To see which second degree polynomials are irreducible, it is enough to consider the canonical form \(a((\mathbf{X} + (b/2a))^2 - (b^2 - 4ac)/4a^2)\) (cf.~VI, p.~26, Prop. 7).

Remark. --- Translating into the canonical trinomial form yields this stronger result: a necessary and sufficient condition for the polynomial \(aX^{2} + bX + c\) over a given ordered field K to have constant sign in K is that \(b^{2} - 4ac < 0\) , and then the sign of the polynomial is that of a.

\subsection{Vector spaces over an ordered field}

Let K be an ordered field, and let E be a vector space over K. The relation « there exists \(\lambda > 0\) in K such that \(y = \lambda x \gg\) between two elements x and y in the set \(E - \{0\}\) is an equivalence relation. The equivalence classes under this relation are called open half-lines with origin 0; the union of an open half-line and \(\{0\}\) is called a closed half-line (or sometimes simply a half-line) with origin 0. Every vector \(a \neq 0\) contained in an open (resp. closed) half-line \(\Delta\) is called a direction vector of \(\Delta\) , and \(\Delta\) is the set of vectors \(\lambda a\) for all scalars \(\lambda > 0\) (resp. \(\lambda \geqslant 0\) ). Every line D through 0 contains exactly two open (resp. closed) half-lines with origin 0; if \(\Delta\) is one of these, then \(-\Delta\) is the other (called the opposite of \(\Delta\) ).

Now if F is an affine space over K, and E the space of translations of F, then any subset of F of the form \(\Delta = a + \Delta_{0}\) , where \(\Delta_{0}\) is an open (resp. closed) half-line of E, is called an open (resp. closed) half-line with origin \(a \in F\) . The half-line \(\Delta_{0}\) is completely determined by \(\Delta\) (for it is the half-line with direction vector b - a, for any \(b \neq a\) in \(\Delta\) ), and is called the direction of \(\Delta\) ; a direction vector of \(\Delta_{0}\) is also called a direction vector of \(\Delta\) .

Suppose now that E has finite dimension n over K; then it is known (III, p.~518, Cor. 1) that the n-th exterior power \(\bigwedge^{n}E\) is a vector space of dimension 1 over K, and so the union of two opposite closed half-lines of origin 0. These half-lines are called orientations of E; the space E together with a given one of these half-lines \(\Delta\) is said to be oriented; an \(n\) -vector \(z\) is then called positive (resp. negative) under this orientation if it belongs to \(\Delta\) (resp. to \(-\Delta\) ); it is negative (resp. positive) under the opposite orientation.

An orientation of an affine space F over K is by definition an orientation of the space of translations of F; the space F, together with such an orientation, is called an oriented affine space.

Let \(\mathbf{E}\) be an oriented vector space over \(\mathbf{K}\) , of dimension \(n\) ; an ordered basis \((a_i)_{1 \leqslant i \leqslant n}\) of \(\mathbf{E}\) is called positive or direct (resp. negative or inverse) if the \(n\) -vector \(a_1 \wedge a_2 \wedge \ldots \wedge a_n\) is positive (resp. negative). If \(u\) is an automorphism of the vector space \(\mathbf{E}\) then \((\bigwedge^{n} u)(z) = \det(u) \cdot z\) for all \(z \in \bigwedge^{n} \mathbf{E}\) , so a necessary and sufficient

condition for \(\bigwedge^{n} u\) to leave invariant the orientation of E (or, as we also say, to preserve the orientation) is that \(\det(u) > 0\) ; the automorphisms having this property are precisely the automorphisms of E as an oriented vector space; they form a normal subgroup \(\mathbf{GL}^{+}(\mathbf{E})\) of the linear group \(\mathbf{GL}(\mathbf{E})\) , which has index 2 whenever \(E \neq 0\) .

When \(\mathbf{E} = 0\) then \(\mathbf{GL}(\mathbf{E}) = \operatorname{End}(\mathbf{E})\) contains only the identity map \(1_{\mathbf{E}}\) and by definition \(\det(1_{\mathbf{E}}) = 1\) . Note that \(\bigwedge^{n}\mathbf{E} = \bigwedge^{0}\mathbf{E} = \mathbf{K}\) by definition in this case; the half-line of \(\mathbf{K}\) formed by the positive scalars is called the canonical orientation of the zero space.

Let \(\mathbf{M}\) and \(\mathbf{N}\) be two complementary subspaces of dimensions \(p\) and \(n - p\) respectively in the vector space \(\mathbf{E}\) of dimension \(n\) ; if \(z'\) (resp. \(z''\) ) is a nonzero vector in \(\bigwedge^{p} \mathbf{M}\) (resp. \(\bigwedge^{n - p} \mathbf{N}\) ), then \(z' \wedge z''\) is a nonzero vector in

\(\bigwedge^{n} \mathbf{E}\). Given an orientation on M and an orientation on N, the vectors \(z' \wedge z''\) for positive \(z'\) and \(z''\) form an orientation of E, called the product orientation of the orientation of M by the orientation of N (which depends on the order of the factors when \(p(n - p)\) is odd). Conversely, given orientations on E and on M, there exists a unique orientation on N such that the given orientation on E is the product of the given orientation on M and this orientation on N (in that order); this orientation is said to be complementary to the orientation of M with respect to that of E. If \(N'\) is a second complementary subspace to M, the canonical projection \(N \to N'\) parallel to M takes the complementary orientation of N onto that of \(N'\) . The image of the complementary orientation of N under the canonical map \(N \to E/M\) is thus independent of the choice of complement N; it is called the quotient orientation on E/M of the orientation of E by that of M.

\BourbakiExercises

\BourbakiExerciseGroup

\begin{enumerate}
\def\labelenumi{\arabic{enumi})}
\tightlist
\item
  A not necessarily commutative ordered monoid M is a (multiplicatively written) monoid equipped with an ordering such that the relation \(x \leqslant y\) implies \(zx \leqslant zy\) and \(xz \leqslant yz\) for all \(z \in M\) . Let G be a not necessarily commutative ordered group. Show that :
\end{enumerate}

\begin{enumerate}
\def\labelenumi{\alph{enumi})}
\item
  If P denotes the set of elements greater than the identity e of G, then P.P = P, \(P \cap P^{-1} = \{e\}\) and \(aPa^{-1} = P\) for all \(a \in G\) . Prove the converse. Find a condition for G to be totally ordered.
\item
  If one of the elements \(\sup (x,y)\) , \(\inf (x,y)\) exists, then so does the other, and
\end{enumerate}

\[
\sup (x, y) = x (\inf (x, y)) ^ {- 1} y = y (\inf (x, y)) ^ {- 1} x.
\]

\begin{enumerate}
\def\labelenumi{\alph{enumi})}
\setcounter{enumi}{2}
\tightlist
\item
  The elements \(\sup(x, e)\) and \(\sup(x^{-1}, e)\) commute. Two coprime elements commute. d) The subgroup \(G'\) generated by the irreducible elements of G is commutative. Deduce that Th. 2 of VI, p.~18 is still valid.
\end{enumerate}

\begin{enumerate}
\def\labelenumi{\arabic{enumi})}
\setcounter{enumi}{1}
\tightlist
\item
  Let E be a lattice with a composition law \((x, y) \mapsto xy\) (not necessarily associative), such that, for all \(a \in E\) , the maps \(x \mapsto ax\) and \(x \mapsto xa\) are order-automorphisms of E. Let \(x_{a}\) (resp. \(_{a}x\) ) denote the element of E defined by \((x_{a})a = x\) (resp. \(a(_{a}x) = x\) ) and suppose that for each \(a \in E\) the maps \(x \mapsto a_{x}\) and \(x \mapsto_{x} a\) are isomorphisms from the ordering of E to the opposite ordering. Show that, for arbitrary x, y, z,
\end{enumerate}

\[
\left(z _ {\inf (x, y)}\right) \cdot x = \left(z _ {y}\right) \cdot \sup (x, y).
\]

\begin{enumerate}
\def\labelenumi{\arabic{enumi})}
\setcounter{enumi}{2}
\item
  Let \(G\) be an ordered group such that the set \(P\) of positive elements is not 0; show that \(G\) is an infinite group, and cannot have a greatest (or a least) element.
\item
  Let G be an ordered group, let P be the set of positive elements of G, and let f be the natural map from G onto a quotient group G/H. For \(f(P)\) to make G/H an ordered group it is necessary and sufficient that \(0 \leqslant y \leqslant x\) and \(x \in H\) imply \(y \in H\) ; then H is called a convex subgroup of G, and G/H is considered as an ordered group in this way. If G is lattice ordered, then a necessary and sufficient condition for G/H to be lattice ordered and for \(f(\sup(x, y)) = \sup(f(x), f(y))\) is that the relations \(|y| \leqslant |x|\) and \(x \in H\) imply \(y \in H\) ; show that this says that H is a filtered convex subgroup. Show that a lattice ordered group with no filtered convex subgroups other than itself and \(\{0\}\) is totally ordered (consider filtered convex subgroups generated by two positive elements of G); * deduce (Gen.~Top., V, p.~25, Ex. 1) that G is then isomorphic to an additive subgroup of the real numbers. *
\item
  Give an example of an ordered group whose ordering is non-discrete, having nonzero elements of finite order (take the quotient of a suitable ordered group by a subgroup H such that \(P \cap H = \{0\}\) ).
\item
  Let G be an ordered group, let P be its set of positive elements. Show that P - P is the largest filtered subgroup of G, and that it is convex. What is the order relation on the quotient?
\item
  In the group \(\mathbf{Z}\) , if \(\mathbf{P}\) is taken to be the set consisting of 0 and the integers \(\geqslant 2\) , then the resulting ordered group is filtered but not lattice ordered (show that the set of elements \(x\) such that \(x \geqslant 0\) and \(x \geqslant 1\) has two distinct minimal elements).
\item
  Let \(x\) be an element of an ordered group such that \(y = \inf(x, 0)\) is defined; if \(n > 0\) is an integer, then \(nx \geqslant 0\) implies \(x \geqslant 0\) (we have
\end{enumerate}

\[
n y = \inf (n x, (n - 1) x, \dots , 0) \geqslant \inf ((n - 1) x, \dots , 0) = (n - 1) y);
\]

and hence \(nx = 0\) implies \(x = 0\) .

\begin{enumerate}
\def\labelenumi{\arabic{enumi})}
\setcounter{enumi}{8}
\item
  In a lattice ordered group G, show that the sum of a family \((H_{\alpha})\) of filtered convex subgroups is a filtered convex subgroup (use VI, p.~12, Cor.).
\item
  Show that, for every finite sequence \((x_{i})\) ( \(1 \leqslant i \leqslant n\) ) of elements of a lattice ordered group \(G\) ,
\end{enumerate}

\[
\sup \left(x _ {i}\right) = \sum_ {i} x _ {i} - \sum_ {i <   j} \inf \left(x _ {i}, x _ {j}\right) + \dots + (- 1) ^ {p + 1} \sum_ {i _ {1} <   i _ {2} <   \dots <   i _ {p}} \inf \left(x _ {i _ {1}}, \dots , x _ {i _ {p}}\right) +   + \dots + (- 1) ^ {n + 1} \inf \left(x _ {1}, \dots , x _ {n}\right).
\]

(Argue by induction based on Prop. 7 (VI, p.~10), using the distributivity of sup over inf.)

\begin{enumerate}
\def\labelenumi{\arabic{enumi})}
\setcounter{enumi}{10}
\item
  Let \((x_{i})\) be a family of n elements of a lattice ordered group G; for each integer k \((0 \leqslant k \leqslant n)\) , let \(d_{k}\) (resp. \(m_{k}\) ) denote the inf (resp. sup) of the \(\binom{n}{k}\) sums of k terms \(x_{i}\) with distinct indices. Show that \(d_{k} + m_{n-k} = x_{1} + x_{2} + \cdots + x_{n}\) .
\item
  Given a lattice ordered group \(\mathbf{G}\) , a subgroup \(\mathbf{H}\) of \(\mathbf{G}\) is called a colattice if for each pair \(x, y\) of elements of \(\mathbf{H}\) , \(\sup_{\mathbf{G}}(x, y) \in \mathbf{H}\) (and so equals \(\sup_{\mathbf{H}}(x, y)\) ).
\end{enumerate}

\begin{enumerate}
\def\labelenumi{\alph{enumi})}
\item
  If \(\mathbf{G} = \mathbf{Q} \times \mathbf{Q} \times \mathbf{Q}\) (where \(\mathbf{Q}\) has the usual ordering), then the subgroup \(\mathbf{H}\) of \((x, y, z)\) such that \(z = x + y\) is a lattice but not a colattice.
\item
  Every filtered convex (Ex. 4) subgroup H of a lattice ordered group G is a colattice.
\item
  Let G be a lattice ordered group and H an arbitrary subgroup of G. Define H' to be the set of infima of finite subsets of H, and H'' to be the set of suprema of finite subsets of H'. Show that H'' is the smallest colattice of G containing H (use the remark on VI, p.~16).
\end{enumerate}

\begin{enumerate}
\def\labelenumi{\arabic{enumi})}
\setcounter{enumi}{12}
\tightlist
\item
  \begin{enumerate}
  \def\labelenumii{\alph{enumii})}
  \tightlist
  \item
    A monoid \(\mathbf{M}\) is said to be lower semi-lattice ordered (or simply semi-lattice ordered) if it is an ordered monoid, if \(\inf(x, y)\) exists for all \(x, y \in \mathbf{M}\) , and if
  \end{enumerate}
\end{enumerate}

\[
\inf (x + z, y + z) = \inf (x, y) + z
\]

for all \(x, y, z \in \mathbf{M}\) . Prove the identities:

\[
\inf (x, z) + \inf (y, z) = \inf (x + y, z + \inf (x, y, z))
\]

\[
\inf (x, y, z) + \inf (x + y, y + z, z + x) = \inf (x, y) + \inf (y, z) + \inf (z, x).
\]

From the first of these relations, deduce that \(x \leqslant z\) and \(y \leqslant z\) imply \(x + y \leqslant z + \inf(x, y)\) .

Show that Prop. 11 and Cor. (VI, pp.~14-15) are valid in a semi-lattice ordered monoid. b) Let M be a semi-lattice ordered monoid and N a submonoid such that \(\inf_{\mathbf{M}}(x,y) = \inf_{\mathbf{N}}(x,y)\) for \(x,y\in \mathbf{N}\) and such that the elements of N are invertible in M. Show that the subgroup G of M consisting of elements \(x - y\) for \(x,y\in \mathbf{N}\) , is lattice ordered.

\begin{enumerate}
\def\labelenumi{\arabic{enumi})}
\setcounter{enumi}{13}
\tightlist
\item
  Show that, in a semi-lattice ordered monoid M,
\end{enumerate}

\[
\inf \left(x _ {i} + y _ {i}\right) \geqslant \inf \left(x _ {i}\right) + \inf \left(y _ {i}\right)
\]

for all finite sequences of n terms \((x_{i})\) , \((y_{i})\) . Deduce the inequalities

\[
(x + y) ^ {+} \leqslant x ^ {+} + y ^ {+}, \quad | x ^ {+} - y ^ {+} | \leqslant | x - y |
\]

in any lattice ordered group.

\begin{enumerate}
\def\labelenumi{\arabic{enumi})}
\setcounter{enumi}{14}
\tightlist
\item
  Show that
\end{enumerate}

\[
\left| x ^ {+} - y ^ {+} \right| + \left| x ^ {-} - y ^ {-} \right| = | x - y |
\]

in a lattice ordered group (notice that \(x - y \leqslant |x - y|\) and \(|x| - |y| \leqslant |x - y|\) ).

\begin{enumerate}
\def\labelenumi{\arabic{enumi})}
\setcounter{enumi}{15}
\tightlist
\item
  Show that
\end{enumerate}

\[
n. \inf (x, y) + \inf (n x, n y) = 2 n. \inf (x, y)
\]

in a semi-lattice ordered monoid (cf.~Ex. 8). Deduce that \(\inf(nx, ny) = n \cdot \inf(x, y)\) if \(\inf(x, y)\) is a regular element.

\begin{enumerate}
\def\labelenumi{\arabic{enumi})}
\setcounter{enumi}{16}
\tightlist
\item
  Let \(x, y, z, t\) be elements of a semi-lattice ordered monoid \(M\) such that \(z \geqslant 0\) and \(t \geqslant 0\) . Prove the inequality
\end{enumerate}

\[
\inf (x + z, y + t) + \inf (x, y) \geqslant \inf (x + z, y) + \inf (x, y + t).
\]

\begin{enumerate}
\def\labelenumi{\arabic{enumi})}
\setcounter{enumi}{17}
\item
  Let \((\mathbf{G}_{\iota})_{\iota \in I}\) be a family of totally ordered groups indexed by a well-ordered set \(I\) ; define an ordering on the direct sum \(\mathbf{G}'\) of the \(\mathbf{G}_{\iota}\) by taking the strictly positive elements to be the nonzero \((x_{\iota})\) such that \(x_{\iota} > 0\) for the first index \(\iota\) such that \(x_{\iota} \neq 0\) . Show that \(\mathbf{G}'\) is totally ordered.
\item
  Let \(G\) be an additive group and \((P_{\alpha})\) a nonempty family of subsets of \(G\) such that \(P_{\alpha} + P_{\alpha} \subset P_{\alpha}\) and \(P_{\alpha} \cap (-P_{\alpha}) = \{0\}\) ; let \(G_{\alpha}\) denote the ordered group obtained from \(G\) by taking \(P_{\alpha}\) as the set of positive elements. Show that \(P + P \subset P\) and \(P \cap (-P) = \{0\}\) , where \(P = \bigcap_{\alpha} P_{\alpha}\) ; if \(H\) is the ordered group obtained from \(G\) by taking \(P\) as the set of
\end{enumerate}

positive elements, show that H is isomorphic to the diagonal subgroup of the product of the \(G_{\alpha}\) .

\begin{enumerate}
\def\labelenumi{\arabic{enumi})}
\setcounter{enumi}{19}
\tightlist
\item
  Let \(G\) be an additive group and let \(P\) be a subset of \(G\) satisfying the following conditions:
\end{enumerate}

\begin{enumerate}
\def\labelenumi{\arabic{enumi}.}
\tightlist
\item
  \(\mathrm{P} + \mathrm{P} = \mathrm{P}\) ; 2. \(\mathrm{P} \cap (-\mathrm{P}) = \{0\}\) ; 3. for every integer \(n > 0\) , \(nx \in \mathrm{P}\) implies \(x \in \mathrm{P}\) (conditions (C)).
\end{enumerate}

\begin{enumerate}
\def\labelenumi{\alph{enumi})}
\item
  If \(a\) is an element of \(G\) such that \(a \notin P\) , show that there exists a subset \(P'\) of \(G\) , satisfying conditions (C), such that \(P \subset P'\) and \(-a \in P'\) (take \(P'\) to be the set of \(x \in G\) such that there exist integers \(m > 0\) and \(n \geqslant 0\) and an element \(y \in P\) with \(mx = -na + y\) .
\item
  Deduce from a) that P is the intersection of those subsets T of G such that \(T + T = T\) , \(T \cap (-T) = \{0\}\) , \(T \cup (-T) = G\) (that is, making G a totally ordered group) and \(P \subset T\) (use Zorn's Lemma).
\item
  In particular if G is an additive group in which every nonzero element has infinite order, then the intersection of all the subsets T of G such that \(T + T = T\) , \(T \cap (-T) = \{0\}\) and \(T \cup (-T) = G\) is \(\{0\}\) .
\end{enumerate}

\begin{enumerate}
\def\labelenumi{\arabic{enumi})}
\setcounter{enumi}{20}
\item
  An ordered group is called lattice orderable if it is isomorphic to a subgroup of a lattice ordered group. Show that it is necessary and sufficient for an ordered group to be lattice orderable that, for every integer \(n > 0\) , the relation \(nx \geqslant 0\) imply \(x \geqslant 0\) (to show that the condition is sufficient, use Ex. 20, b) and 19). Show that every lattice orderable group is isomorphic to a subgroup of a product of totally ordered groups.
\item
  Let G be a lattice orderable group (considered as a Z-module) and E the vector space \(G_{(Q)}\) (II, p.~277); show that there is a unique ordering on E compatible with the additive group structure of E, inducing the given ordering on G, and such that E is lattice orderable.
\item
  Let G be the lattice ordered group \(Z \times Z\) (where Z has the usual ordering) and H the convex subgroup of G generated by \((2, -3)\) ; show that the ordered group G/H is not lattice orderable (cf.~VI, p.~30, Ex. 4).
\item
  Let A be a commutative ring and I the set of all ideals of A. Then \(\mathfrak{a}(\mathfrak{b} + \mathfrak{c}) = \mathfrak{a}\mathfrak{b} + \mathfrak{a}\mathfrak{c}\) for all a, b, c in I, by I, p.~107; in other words I, with the order relation \(a \supset b\) and the law of composition \((\mathfrak{a}, \mathfrak{b}) \mapsto \mathfrak{a}\mathfrak{b}\) , is a semi-lattice ordered monoid (VI, p.~31, Ex. 13), the sup of two elements a, b of I being \(a \cap b\) , and their inf \(a + b\) .
\end{enumerate}

\begin{enumerate}
\def\labelenumi{\alph{enumi})}
\tightlist
\item
  Let K be a field and \(A = K[X, Y]\) the ring of polynomials in two indeterminates over K.
\end{enumerate}

\[
\mathbf {a} = (\mathbf {X}), \mathbf {b} = (\mathbf {Y}), \mathbf {c} = (\mathbf {X} + \mathbf {Y})
\]

\[
\begin{array}{c} (\mathfrak {a} \cap \mathfrak {b}) + \mathfrak {c} \neq (\mathfrak {a} + \mathfrak {c}) \cap (\mathfrak {b} + \mathfrak {c}), \\ (\mathfrak {a} + \mathfrak {b}) \cap \mathfrak {b} \neq (\mathfrak {a} \cap \mathfrak {c}) + (\mathfrak {b} \cap \mathfrak {c}), \\ (\mathfrak {a} \cap \mathfrak {b}) (\mathfrak {a} + \mathfrak {b}) \neq (\mathfrak {a} (\mathfrak {a} + \mathfrak {b})) \cap (\mathfrak {b} (\mathfrak {a} + \mathfrak {b})). \end{array}
\]

\begin{enumerate}
\def\labelenumi{\alph{enumi})}
\setcounter{enumi}{1}
\tightlist
\item
  In the ring A, 1 is a gcd of X and Y, but \((\mathbf{X}) + (\mathbf{Y}) \neq \mathbf{A}\) .
\end{enumerate}

\begin{enumerate}
\def\labelenumi{\arabic{enumi})}
\setcounter{enumi}{24}
\tightlist
\item
  Let I be the semilattice ordered monoid of ideals of a commutative ring A (Ex. 24). For a finite system of congruences \(x \equiv a_i(a_i)\) to admit a solution whenever any two of them admit a common solution (that is to say whenever
\end{enumerate}

\[
a _ {i} \equiv a _ {j} (\mathfrak {a} _ {i} + \mathfrak {a} _ {j})
\]

for every pair of indices \(i, j)\) ; it is necessary and sufficient that each of the laws

\[
(a, b) \mapsto a \cap b, (a, b) \mapsto a + b,
\]

in I be distributive over the other (« Chinese remainder theorem »). One may proceed as follows :

\begin{enumerate}
\def\labelenumi{\alph{enumi})}
\item
  If \((\mathfrak{a}_1 \cap \mathfrak{a}_2) + (\mathfrak{a}_1 \cap \mathfrak{a}_3) = \mathfrak{a}_1 \cap (\mathfrak{a}_2 + \mathfrak{a}_3)\) and if any two of the three congruences \(x \equiv a_i(\mathfrak{a}_i)\) ( \(i = 1, 2, 3\) ) admit a common solution, then the three congruences have a common solution (let \(x_{12}\) be a common solution of \(x \equiv a_1(\mathfrak{a}_1)\) and \(x \equiv a_2(\mathfrak{a}_2)\) and let \(x_{13}\) be a common solution of \(x \equiv a_1(\mathfrak{a}_i)\) and \(x \equiv a_3(\mathfrak{a}_3)\) ; show that the congruences \(x \equiv x_{12}(\mathfrak{a}_1 \cap \mathfrak{a}_2)\) and \(x \equiv x_{13}(\mathfrak{a}_1 \cap \mathfrak{a}_3)\) have a common solution).
\item
  If any system of three congruences, any two of which have a common solution, admits a common solution, show that the following distributive laws hold
\end{enumerate}

\[
\mathfrak {a} + (\mathfrak {b} \cap \mathfrak {c}) = (\mathfrak {a} + \mathfrak {b}) \cap (\mathfrak {a} + \mathfrak {c}), \quad \text { and } \quad \mathfrak {a} \cap (\mathfrak {b} + \mathfrak {c}) = (\mathfrak {a} \cap \mathfrak {b}) + (\mathfrak {a} \cap \mathfrak {c}).
\]

(For the first, note that for all \(x \in (\mathfrak{a} + \mathfrak{b}) \cap (\mathfrak{a} + \mathfrak{c})\) , there exists \(y \in \mathfrak{b} \cap \mathfrak{c}\) such that \(y \equiv x(\mathfrak{a})\) ; for the second, note that for all \(x \in \mathfrak{a} \cap (\mathfrak{b} + \mathfrak{c})\) , there exists \(y \in \mathfrak{a} \cap \mathfrak{b}\) such that \(y \equiv x(\mathfrak{c})\) .) Note that by \(a)\) and \(b)\) , the second distributive law implies the first (cf.~Set Theory, III, p.~217, Ex. 16).

\begin{enumerate}
\def\labelenumi{\alph{enumi})}
\setcounter{enumi}{2}
\item
  Prove the « Chinese remainder theorem » by induction on the number of congruences under consideration, by an analogous argument to that of a).
\item
  * Translation in the case of principal ideal domains. *
\end{enumerate}

\begin{enumerate}
\def\labelenumi{\arabic{enumi})}
\setcounter{enumi}{25}
\item
  Show that the ideal (X) in the monoid of ideals of the polynomial ring K{[}X, Y{]} (where K is a commutative field) satisfies condition (P) of Prop. 14 of VI, p.~17, but is not maximal.
\item
  Let A be the quadratic Z-algebra with basis \((1, e)\) , where \(e^{2} = -5\) . Show that A is an integral domain; in A we have \(9 = 3 \cdot 3 = (2 + e)(2 - e)\) ; show that 3, \((2 + e)\) and \((2 - e)\) are irreducible elements of A, but do not satisfy the condition in Prop. 14 (DIV) of VI, p.~17.
\item
  Let G be a lattice ordered group with positive elements P. Two elements x and y of P are said to be equivalent if any element coprime to one is also coprime to the other; the equivalence classes under this relation are called the threads of P; let \(\bar{x}\) denote the thread containing x.
\end{enumerate}

\begin{enumerate}
\def\labelenumi{\alph{enumi})}
\item
  Let \(\overline{a}\) and \(\overline{b}\) be threads; let \(x, x_1\) be elements of \(\overline{a}\) and \(y, y_1\) be elements of \(\overline{b}\) ; show that if every element coprime to \(x\) is also coprime to \(y\) , then every element coprime to \(x_1\) is also coprime to \(y_1\) . Define a relation in this way between \(\overline{a}\) and \(\overline{b}\) , written \(\overline{a} \geqslant \overline{b}\) ; show that this is an ordering of the set F of threads.
\item
  Show that \(\mathbf{F}\) is a lattice under this ordering, and that
\end{enumerate}

\[
\inf (\bar {a}, \bar {b}) = \overline {{\inf (a , b)}}, \quad \sup (\bar {a}, \bar {b}) = \overline {{\sup (a , b)}} = \overline {{a + b}}
\]

(use Cor. 3 of VI, p.~15). Show that F has a smallest element, namely the thread \(\overline{0} = \{0\}\) .

\begin{enumerate}
\def\labelenumi{\alph{enumi})}
\setcounter{enumi}{2}
\item
  Call two threads \(\overline{a}\) , \(\overline{b}\) coprime if \(\inf (\overline{a},\overline{b}) = \overline{0}\) . Show that if \(\overline{a}\) and \(\overline{b}\) are two threads \(\neq \overline{0}\) , and if \(\overline{a} < \overline{b}\) , then there exists a thread \(\overline{c}\) , coprime to \(\overline{a}\) , such that \(\overline{c} < \overline{b}\) .
\item
  Call a thread \(\bar{m} \neq \bar{0}\) irreducible if it is a minimal element of the set of threads \(\neq \bar{0}\) . Show that an irreducible thread is totally ordered (if \(a\) and \(b\) are two elements of \(\bar{m}\) , consider the threads containing \(a - \inf(a, b)\) and \(b - \inf(a, b)\) and use \(b\) ). Show also that the union of \(\bar{m}, -\bar{m}\) and \(\{0\}\) is a totally ordered subgroup \(\mathrm{H}(\bar{m})\) of \(\mathbf{G}\) .
\item
  Given a family \((\bar{m}_{\iota})\) of distinct irreducible threads of G, show that the subgroup H generated by the union of the \(\bar{m}_{\iota}\) is isomorphic to the direct sum of the groups \(\mathrm{H}(\bar{m}_{\iota})\) (reduce to the case of a finite family, and argue by induction).
\item
  Show that in a product (resp. direct sum) of nonzero totally ordered groups, the elements of irreducible threads are those in which all the coordinates are zero except for one (which is positive). Deduce that an ordered group can be a product (resp. direct sum) of nonzero totally ordered groups in at most one way (in other words the factors are uniquely determined).
\end{enumerate}

\begin{enumerate}
\def\labelenumi{\arabic{enumi})}
\setcounter{enumi}{28}
\tightlist
\item
  An ordered group G is completely lattice ordered if it is filtered and if every nonempty subset of G which is bounded above has a supremum in G. A lower semi-lattice ordered monoid is said to be completely semi-lattice ordered if \(\inf(x_{\alpha})\) exists for every nonempty
\end{enumerate}

family of elements \((x_{\alpha})\) of M which is bounded below, and if

\[
\inf _ {\alpha , \beta} (x _ {\alpha} + y _ {\beta}) = \inf _ {\alpha} (x _ {\alpha}) + \inf _ {\beta} (y _ {\beta})
\]

for any two nonempty families \((x_{\alpha})\) , \((y_{\beta})\) which are bounded below. A completely lattice ordered group is a completely semi-lattice ordered monoid.

\begin{enumerate}
\def\labelenumi{\alph{enumi})}
\item
  Show that a necessary and sufficient condition for a filtered group G to be completely lattice ordered is that every nonempty subset of the set P of positive elements of G admit an infimum in P.
\item
  Every product of completely lattice ordered groups is completely lattice ordered.
\item
  Every filtered convex subgroup of a completely lattice ordered group is completely lattice ordered.
\end{enumerate}

\begin{enumerate}
\def\labelenumi{\arabic{enumi})}
\setcounter{enumi}{29}
\tightlist
\item
  Let G be an ordered group. For every subset A of G let \(m(A)\) (resp. M(A)) denote the set of lower (resp. upper) bounds for A in G; then \(m(A) = -M(-A)\) . a) The relation \(A \subset B\) implies \(m(B) \subset m(A)\) and \(\mathbf{M}(m(A)) \subset \mathbf{M}(m(B))\) .
\end{enumerate}

\begin{enumerate}
\def\labelenumi{\alph{enumi})}
\setcounter{enumi}{1}
\item
  We have \(\mathbf{A} \subset \mathbf{M}(m(\mathbf{A}))\) and \(\mathbf{M}(m(\mathbf{M}(\mathbf{A}))) = \mathbf{M}(\mathbf{A})\) .
\item
  If \((\mathbf{A}_{\iota})\) is an arbitrary family of subsets of \(G\) then \(m\left(\bigcup_{\iota} \mathbf{A}_{\iota}\right) = \bigcap_{\iota} m(\mathbf{A}_{\iota})\) .
\item
  Every set M(B), where B is a nonempty subset bounded above in G, is called a major set in G. For any nonempty subset A bounded below in G, the set M(m(A)) is the smallest major set containing A; denote it \(\langle A\rangle\) ; if \(A\subset B\) then \(\langle A\rangle\subset\langle B\rangle\) ; if A has an infimum a in G, then \(\langle A\rangle=M(a)\) , which we also write \(\langle a\rangle\) .
\item
  If A and B are two nonempty subsets bounded below in G, then \(\langle A + B \rangle = \langle \langle A \rangle + B \rangle\) . Deduce that, in the set \(\mathfrak{M}(G)\) of major sets in G, the map \((A, B) \mapsto \langle A + B \rangle\) is an associative and commutative composition law, for which \(\langle 0 \rangle = P\) is an identity, that the order relation \(A \supset B\) is compatible with this law, and that \(\mathfrak{M}(G)\) is a completely semi-lattice ordered monoid (Ex. 29) under this structure if G is filtered. Moreover, the map \(x \mapsto \langle x \rangle\) from G to \(\mathfrak{M}(G)\) is an isomorphism from the ordered group G onto a subgroup of the monoid \(\mathfrak{M}(G)\) .
\item
  If an element A of \(\mathfrak{M}(G)\) is invertible under the composition law of this monoid, then its inverse is \(\mathbf{M}(-\mathbf{A}) = -m(\mathbf{A})\) (notice that if B is the inverse of A, then the relations \(\mathbf{B} \subset \mathbf{M}(-\mathbf{A})\) and \(\mathbf{A} + \mathbf{M}(-\mathbf{A}) \subset \langle 0 \rangle\) hold, whence \(\langle \mathbf{A} + \mathbf{B} \rangle = \langle \mathbf{A} + \mathbf{M}(-\mathbf{A}) \rangle\) ).
\item
  For an element A of \(\mathfrak{M}(G)\) to be invertible, it is necessary and sufficient that \(x + A \subset A\) imply \(x \geqslant 0\) (write down an expression for \(0 \in \langle A + M(-A) \rangle\) ).
\end{enumerate}

\begin{enumerate}
\def\labelenumi{\arabic{enumi})}
\setcounter{enumi}{30}
\tightlist
\item
  An ordered group G is called archimedean if the only elements \(x \in G\) such that the set of nx (n a positive integer) is bounded below are the positive elements of G.
\end{enumerate}

\begin{enumerate}
\def\labelenumi{\alph{enumi})}
\item
  The monoid \(\mathfrak{M}(G)\) of major sets in an ordered group G is a group if and only if G is archimedean (use Ex. 30, g) and d). If in addition G is filtered, then \(\mathfrak{M}(G)\) is a completely lattice ordered group.
\item
  Deduce that a filtered group G is isomorphic to a subgroup of a completely lattice ordered group if and only if G is archimedean.
\end{enumerate}

\begin{enumerate}
\def\labelenumi{\arabic{enumi})}
\setcounter{enumi}{31}
\tightlist
\item
  \begin{enumerate}
  \def\labelenumii{\alph{enumii})}
  \tightlist
  \item
    Let G be a completely lattice ordered group, and H a subgroup of G. For every major subset A of the group H, let \(x_{\mathbf{A}}\) denote the infimum of A in G; show that the map \(\mathbf{A} \mapsto x_{\mathbf{A}}\) from \(\mathfrak{M}(\mathbf{H})\) to G is injective (show that A is the set of \(y \in \mathbf{H}\) such that \(y \geqslant x_{\mathbf{A}}\) ).
  \end{enumerate}
\end{enumerate}

\begin{enumerate}
\def\labelenumi{\alph{enumi})}
\setcounter{enumi}{1}
\item
  For every subset B of H, bounded below in H, let \(\langle B\rangle\) denote the major set generated by B in H (an element of \(\mathfrak{M}(H)\) ). If for any subset B of H, bounded below in H, the equation \(\inf B = \inf \langle B \rangle\) holds (infima taken in G), show that the map \(A \mapsto x_{A}\) is an isomorphism of the ordered group \(\mathfrak{M}(H)\) onto a subgroup of G (cf.~Ex. 30, e)).
\item
  If every element of G is the infimum of some subset of H, show that for any subset B of H, bounded below in H, we have inf B = inf \(\langle B\rangle\) (infima taken in G). (Notice that this holds if inf B ∈ H; in general, consider an element \(x \in G\) such that \(x + inf B \in H\) and use the fact that x is the infimum of some subset of G, as well as Ex. 30, e).
\end{enumerate}

\(d) ^{*}\) Let \(G\) be the completely lattice ordered group \(\mathbf{Z} \times \mathbf{Z} \times \mathbf{R}\) . Let \(\theta > 0\) be an irrational number; let \(H\) be the subgroup of \((x, y, z)\) such that \(\theta(z - x) + y = 0\) . Show that there exists no isomorphism of the ordered group \(\mathfrak{M}(H)\) onto any subgroup of \(G\) which restricts to the identity map on \(H\) (notice that \(\mathfrak{M}(H)\) is isomorphic to \(K = Z \times R\) ; consider the subgroup of \(K\) consisting of elements \(u\) such that, for every integer \(n > 0\) , there exists \(v \in K\) with \(nv = u\) , and the analogous subgroup of \(G\) ). *

\begin{enumerate}
\def\labelenumi{\arabic{enumi})}
\setcounter{enumi}{32}
\tightlist
\item
  \begin{enumerate}
  \def\labelenumii{\alph{enumii})}
  \tightlist
  \item
    A totally ordered group G is archimedean if and only if, for any pair \(x > 0\) , \(y > 0\) of elements of G, there exists an integer \(n > 0\) such that \(y \leqslant nx\) .
  \end{enumerate}
\end{enumerate}

* b) Any totally ordered, archimedean, completely lattice ordered group G is isomorphic to \(\{0\}\) , Z or R (ruling out the first two cases, fix a \textgreater{} 0 in G, and make each \(x \in G\) correspond to the infimum of the rational numbers p/q such that \(qx \leq pa\) ).

\begin{enumerate}
\def\labelenumi{\alph{enumi})}
\setcounter{enumi}{2}
\item
  Deduce that every archimedean totally ordered group G is isomorphic to a subgroup of R (notice that \(\mathfrak{M}(G)\) is totally ordered). *
\item
  The lexicographic product \(Z \times Z\) is totally ordered and not archimedean.
\end{enumerate}

\begin{enumerate}
\def\labelenumi{\arabic{enumi})}
\setcounter{enumi}{33}
\item
  Let \(G = Z^{N}\) be the product of a countably infinite family of totally ordered groups Z, and let \(H = Z^{(N)}\) be the filtered convex subgroup, the direct sum of the factors of G. Show that the lattice ordered group G/H (VI, p.~30, Ex. 4) is not archimedean, although G and H are completely lattice ordered.
\item
  A major set A in an ordered group G is said to be finitely generated if there exists a finite set F such that \(A = \langle F \rangle\) . We say that G is semi-archimedean if every finitely generated major set is invertible in the monoid \(\mathfrak{M}(G)\) . Every lattice ordered (resp. archimedean) group is semi-archimedean (Ex. 31, a)).
\end{enumerate}

\begin{enumerate}
\def\labelenumi{\alph{enumi})}
\item
  Show that every filtered semi-archimedean group is lattice ordered (if \(nx \geqslant 0\) , consider the major set \(\langle F \rangle\) , where \(F = \{0, x, \ldots, (n - 1)x\}\) , and use Ex. 30, g) and the fact that it is invertible).
\item
  * Let K be the lexicographic product \(R \times R\) , and G the (usual) product group \(K \times R\) . Let \(\theta\) be an irrational number such that \(0 < \theta < 1\) , and let H be the subgroup generated by \(((1, 0), 0)\) , \(((\theta, 0), \theta)\) and \(((0, x), 0)\) where x runs through R. Show that H, a subgroup of a product of two totally ordered groups, is not semi-archimedean (consider the major set generated by the two elements \(((1,0),0)\) and \(((\theta ,0),\theta)\) of H, and show that it is not invertible). *
\item
  Let G be a semi-archimedean filtered group; show that the submonoid of \(\mathfrak{M}(G)\) generated by the finitely generated major sets and their inverses is a lattice ordered group (cf.~VI, p.~31, Ex. 13, b)).
\item
  * Let \(\theta > 0\) be an irrational number, and let \(G\) be the group \(\mathbf{Z} \times \mathbf{Z}\) in which the set \(P\) of positive elements is taken to be the set of \((x, y)\) such that \(x \geqslant 0\) and \(y \geqslant \theta x\) . Show that \(G\) is an archimedean group, but that the inverse in \(\mathfrak{M}(G)\) of a finitely generated major set, not of the form \(\langle a \rangle\) , is not finitely generated. *
\end{enumerate}

\begin{enumerate}
\def\labelenumi{\arabic{enumi})}
\setcounter{enumi}{35}
\tightlist
\item
  Let \(\mathbf{G}\) be a filtered group with positive elements \(\mathbf{P}\) . Then \(\mathbf{G}\) is isomorphic to a group of the form \(\mathbf{Z}^{(I)}\) if and only if there exists a map \(x \mapsto d(x)\) from \(\mathbf{P}\) into \(\mathbf{N}\) such that, for all \(b\) , either \(b \geqslant a\) or there exists an element \(c\) in the major set \(\langle a, b \rangle\) such that \(d(c) < d(a)\) . (Show that this condition is satisfied in \(\mathbf{Z}^{(I)}\) ) by the map \(d(x_{i}) = \sum_{i} x_{i}\) (when \(x_{i} \geqslant 0\) for all \(i\) ).
\end{enumerate}

Conversely, if the condition is satisfied, show that every major set \(\mathbf{A} \subset \mathbf{P}\) has the form \(\langle a \rangle\) , by taking an element \(a\) of \(\mathbf{A}\) such that \(d(a) \leqslant d(x)\) for all \(x \in \mathbf{A}\) . Apply Th. 2 of VI, p.~18.)
